% Part: normal-modal-logic
% Chapter: completeness
% Section: lindenbaums-lemma

\documentclass[../../../include/open-logic-section]{subfiles}

\begin{document}

\olfileid{nml}{com}{lin}

\olsection{مبرهنة ليندنباوم}

نحتاج إلى نقل كل مجموعة $\Sigma$-متسقة من ال!!{formula}s إلى مجموعة
$\Sigma$-متسقة \emph{تامة} تحتويها. وهذا ما تكفله مبرهنة ليندنباوم.
وسنرى عند بناء النموذج القانوني أن لكل مجموعة $\Sigma$-متسقة تامة
$\Delta$ عالمًا تصدق عنده ال!!{formula}s المنتمية إلى~$\Delta$ دون غيرها.
فإذا ثبت التوسيع، ثبت أن كل مجموعة $\Sigma$-متسقة
تصدق عند عالم من عوالم ذلك النموذج.

\begin{thm}[مبرهنة ليندنباوم]\ollabel{thm:lindenbaum}
  إذا كانت~$\Gamma$ $\Sigma$-متسقة، فثمة مجموعة $\Sigma$-متسقة تامة
  $\Delta$ توسّع~$\Gamma$.
\end{thm}

\begin{proof}
لتكن~$!A_0$، $!A_1$، \dots{} قائمة شاملة لجميع صيغ اللغة (ويجوز
التكرار). فمثلًا، نبدأ بإدراج~$\Obj p_0$، وفي كل مرحلة~$n \ge 1$
ندرج الصيغ المنتهية الكثيرة التي لا يتجاوز طولها~$n$ ولا تستعمل إلا
المتغيرات~$\Obj p_0$، \dots،~$\Obj p_n$. ونعرّف مجموعات من
ال!!{formula}s~$\Delta_n$ بالاستقراء على~$n$، ثم نضع
$\Delta = \bigcup_n \Delta_n$. نضع أولًا~$\Delta_0 = \Gamma$.
وبافتراض أن~$\Delta_n$ قد عُرّفت، نعرّف~$\Delta_{n+1}$ هكذا:
\[
\Delta_{n+1} =
\begin{cases}
  \Delta_n \cup \{!A_n\}, & \text{إذا كانت $\Delta_n \cup \{!A_n\}$
    $\Sigma$-متسقة؛} \\
  \Delta_n \cup \{\lnot !A_n\}, & \text{فيما عدا ذلك.}
\end{cases}
\]
والآن لتكن~$\Delta = \bigcup_{n=0}^\infty \Delta_n$.

ولنحقق الآن أن المجموعة~$\Delta$ المعرّفة بهذا البناء تستوفي المطلوب:
أن يكون~$\Gamma \subseteq \Delta$، وأن تكون~$\Delta$ $\Sigma$-متسقة تامة.

أما~$\Gamma \subseteq \Delta$ فظاهر، إذ~$\Delta_0 \subseteq
\Delta$ بالبناء، و$\Delta_0 = \Gamma$.
بل~$\Delta_n \subseteq \Delta$ لكل~$n$، لأن~$\Delta$ اتحاد المجموعات~$\Delta_n$.
والمجموعات في هذا البناء متزايدة؛ فإننا نضيف في كل مرحلة !!a{formula}
إلى ما سبق، فيكون $\Delta_n \subseteq \Delta_{n+1}$،
وبتعدي~$\subseteq$ يكون $\Delta_n \subseteq \Delta_m$ إذا كان~$n \le m$.
وأما التمام، فلأن كل مرحلة تضم إما~$!A_n$ وإما~$\lnot !A_n$،
وقائمتنا لا تخلو من أي !!{formula}، وإن تكرر بعضها في~$!A_n$.
فلكل~$!A$ يكون $!A \in \Delta$ أو~$\lnot !A \in \Delta$؛
وذلك تمام~$\Delta$ بعينه.

وأما اتساق~$\Delta$ بالنسبة إلى~$\Sigma$ فبرهانه من وجهين متعاقبين.
نبين أولًا أن عدم اتساق~$\Delta$ مع~$\Sigma$ يرجع إلى عدم اتساق
إحدى المراحل~$\Delta_n$ مع~$\Sigma$؛ ثم نثبت أن كل~$\Delta_n$
$\Sigma$-متسقة، فيمتنع عدم اتساق الاتحاد.

لنفترض أن~$\Delta$ $\Sigma$-غير متسقة.
إذن $\Delta \Proves[\Sigma] \lfalse$، أي توجد مقدمات متناهية
$!A_1$، \dots،~$!A_k \in \Delta$ تحقق
$\Sigma \Proves !A_1 \lif (!A_2 \lif \cdots (!A_k \lif
\lfalse)\dots)$. وبما أن~$\Delta = \bigcup_{n=0}^\infty \Delta_n$،
فكل~$!A_i \in \Delta_{n_i}$ لمرحلة ما~$n_i$.
نأخذ~$n$ أكبر الفهارس التي احتجنا إليها؛ فإن خلا الاشتقاق من المقدمات
كفت المرحلة الأولى. ولكون~$n_i \le n$ يلزم~$\Delta_{n_i} \subseteq \Delta_n$.
فتجتمع جميع~$!A_i$ في~$\Delta_n$ واحدة،
ويشهد الاشتقاق نفسه على $\Delta_n \Proves[\Sigma] \lfalse$.
فتكون~$\Delta_n$ $\Sigma$-غير متسقة، كما أردنا.

ثم نثبت اتساق كل~$\Delta_n$ بالنسبة إلى~$\Sigma$ بالاستقراء على~$n$.
أما البداية فـ~$\Delta_0 = \Gamma$، وقد فرضنا اتساق~$\Gamma$
بالنسبة إلى~$\Sigma$، فتم المطلوب عند~$n = 0$.
وافرض ثبوته عند~$n$؛ أي اتساق~$\Delta_n$ بالنسبة إلى~$\Sigma$.
بالتعريف، تكون~$\Delta_{n+1}$ هي
$\Delta_n \cup \{!A_n\}$ إن كانت هذه $\Sigma$-متسقة،
وتكون $\Delta_n \cup \{\lnot!A_n\}$ فيما عدا ذلك.
ففي الحال الأولى ثبت اتساق~$\Delta_{n+1}$ بالنسبة إلى~$\Sigma$
بنفس شرط الاختيار. وفي الأخرى نستعمل
\olref[prf][con]{prop:consistencyfacts}\olref[prf][con]{prop:consistencyfacts-c}:
إحدى المجموعتين $\Delta_n \cup \{!A_n\}$ و$\Delta_n \cup \{\lnot!A_n\}$
متسقة؛ فإذا امتنع اتساق الأولى تعيّن اتساق الثانية.
وعليه فـ~$\Delta_{n+1}$ متسقة في الحالين.
\end{proof}

\begin{cor}\ollabel{cor:provability-characterization}
  $\Gamma \Proves[\Sigma] !A$ إذا وفقط إذا كان~$!A \in \Delta$ لكل
  مجموعة $\Sigma$-متسقة تامة~$\Delta$ توسع~$\Gamma$ (بما في ذلك حالة
  $\Gamma = \emptyset$، فنحصل عندئذ على توصيف آخر للنسق الجهوي
  $\Sigma$).
\end{cor}

\begin{proof}
  إذا كان~$\Gamma \Proves[\Sigma] !A$، وأخذنا~$\Delta$ مجموعة
  $\Sigma$-متسقة تامة توسّع~$\Gamma$، امتنع أن يكون~$!A \notin \Delta$.
  إذ يلزم حينئذ~$\lnot!A \in \Delta$ بالقصوية، ولدينا
  $\Delta \Proves[\Sigma] !A$ بالرتابة، و
  $\Delta \Proves[\Sigma] \lnot!A$ بالانعكاسية؛ فتبطل دعوى اتساق~$\Delta$.
  وأما الاتجاه الآخر فنثبته بالمقابل. إذا كان~$\Gamma \Proves/[\Sigma] !A$،
  كانت $\Gamma \cup \{\lnot!A\}$ $\Sigma$-متسقة.
  فبمبرهنة ليندنباوم لها توسعة متسقة تامة~$\Delta$
  تحتوي~$\Gamma \cup \{\lnot!A\}$؛ وباتساقها يكون $!A \notin \Delta$.
\end{proof}

\end{document}
